\section{열린 경계와 과장 금지선}
\label{sec:boundaries}

\subsection{가장 좁지만 중요한 의미론(semantics) 간극(gap): product/replay}

현재 raw \(\mu\) 선에는 다음 두 사실이 각각 존재한다.

\begin{enumerate}
  \item 실제 generated Sign/Verify hash 호출의 \emph{반환값}은 관련 입력에서
        같은 32바이트 접두부를 갖는다;
  \item 실제 Sign 후 Verify 전체 실행은 ghost trace와 exact하게 대응하고,
        trace는 각 hash 결과를 \identifier{observed_mu}에 저장한다.
\end{enumerate}

원하는 결론은 accepted Verify 사후상태(post-state)에서
\[
\begin{aligned}
  \mathit{VerifyResult}=0\Longrightarrow
  \operatorname{MuPrefix}_{32}(&
    \mathit{SignTrace.observed\_mu},\\[-0.3em]
    &\mathit{VerifyTrace.observed\_mu})
\end{aligned}
\]
이다. 수학적으로는 1과 2를 합치면 자명해 보이지만, 현재 program-logic
표면에서는 두 명제(proposition)의 변수(variable)가 같은 시점에 존재하지 않는다.

\begin{center}
\begin{tikzpicture}[node distance=8mm]
  \node[provennode] (returns) {hash-call 시점의\\두 반환값 pRHL 관계};
  \node[partialnode, below=of returns] (stored) {완료된 두 trace의\\stored observation 관계};
  \node[provennode, left=12mm of stored] (signcont) {Sign의 잔여 실행\\signature core};
  \node[provennode, right=12mm of stored] (verifycont) {Verify의 잔여 실행\\challenge와 mismatch};
  \draw[openarrow] (returns) -- node[right,align=left]{product/replay\\운송 부재} (stored);
  \draw[flowarrow] (returns) -- (signcont);
  \draw[flowarrow] (returns) -- (verifycont);
\end{tikzpicture}
\end{center}

hash-call의 pRHL goal에는 \(res\{1\},res\{2\}\)가 있지만 전체 trace가 끝난
unary post-state에는 이 지역 반환변수가 사라진다. 반대로 hash-call 지점에서
pRHL 정리(theorem)를 적용하면 Sign과 Verify의 서로 다른 잔여 continuation을 처리해야
한다. Sign 쪽 continuation을 일방적으로 버리려면 아직 없는 losslessness 또는
적절한 product rule이 필요하다.

\begin{openobligation}[\claimid{OBL-MU-PRODUCT-REPLAY}]
반환값 수준의 region-local pRHL 정리(theorem)를 이미 완료된 sequential trace의 stored
observation 관계(relation)로 운반하는 sound rule 또는 증명 구조가 필요하다.
\end{openobligation}

이 의무는 Week~7에 \statusdeferred / explicit paper boundary로 동결되었다
\cite{product-replay}. 다음 접근은 의도적으로 거부되었다.

\begin{itemize}
  \item observation equality를 전제로 가정하여 결론을 재진술하기;
  \item SHAKE locality axiom을 새로 추가하기;
  \item 또 다른 caller/observation wrapper로 theorem surface를 계속 넓히기;
  \item termination 근거 없이 Sign continuation을 버리기.
\end{itemize}

따라서 \(\delta_{\mu,\mathrm{raw\_api\_accept}}=0\)은 현재 주장하지 않는다.
이것은 byte-locality 정리(theorem)가 부족해서가 아니라 프로그램 논리의 운송 화살표가
없어서 생긴 의미론 간극이다.

\subsection{accepted-path와 legitimate-signature correctness는 다른 방향이다}

실제 Verify에 대해 증명된 방향은
\[
  \operatorname{Accept}\Longrightarrow\operatorname{TailReached}
  \Longrightarrow\operatorname{HashInputsBound}
\]
이다. 이는 이미 accept한 실행을 분석하는 security-facing lane이다.

기능정확성에 필요한 방향은
\[
  \sigma\leftarrow\operatorname{Sign}(sk,M)
  \Longrightarrow
  \operatorname{Verify}(vk,M,\sigma)=\operatorname{Accept}
\]
이다. 이를 위해서는 Sign이 만든 서명이 unpack/norm/hint/challenge gate를 모두
통과한다는 사실이 필요하다. 첫 방향의 정리(theorem)를 둘째 방향으로 뒤집을 수 없다.

\begin{openobligation}[\claimid{OBL-SIGN-OUTPUT-TAIL-REACH}]
실제 Sign output이 실제 Verify의 early-reject tree를 통과해 tail에 도달함을
증명해야 한다. 현재 signature prefix inverse와 성공 조건부 full-HBZ wrapper는
닫혔고, fixed all-zero HBZ의 반환 실행에서는 failure도 배제되었다. 그러나 이
고정 입력은 실제 Sign output 전체의 canonicality를 대신하지 않는다. \(h\)
suffix, metadata/padding, norm, hint, arithmetic correctness가 남아 있다.
\end{openobligation}

\claimid{OBL-SIGN-VERIFY-CORRECTNESS}는 이 의무와 나머지 산술·challenge
정확성에 막혀 \statusblocked 다.

\subsection{full signature codec에서 남은 표현론}

prefix의 단면--수축(section--retraction)을 full codec으로 확장하려면 적어도
다음 사각형(square)이 필요하다. Week~14에는 이 중 HBZ 간선 하나가 닫혔고,
Week~15에는 그 간선의 fixed all-zero 반환 실행에서 failure가 배제되었다.

\begin{itemize}
  \item 두 size-delta metadata byte가 pack/unpack에서 같은 크기를 나타냄;
  \item HBZ rANS actual full wrapper가 성공 조건 아래 역함수(inverse function)임
        (\statusproved, Week~14);
  \item canonical all-zero HBZ의 fixed-input Hoare success
        (\statusproved, Week~15; termination은 미증명);
  \item 두 번째 \(h\) suffix codec의 encode/decode
        (\statusdeferred, Week~16 MINCORE 전환);
  \item trailing zero padding의 정준성(canonicality)과 parser의 zero gate;
  \item packer/unpacker가 목표 밖 배열 tail을 보존함;
  \item 실제 Sign이 canonical 입력과 성공 buffer bound를 산출함.
\end{itemize}

따라서
\claimid{OBL-SIG-PREFIX-CODEC}은 제한된 의미에서 \statusproved 지만,
\claimid{OBL-SIG-FULL-CODEC-MODE2}는 \statusblocked 다. prefix inverse와 full
inverse를 같은 명제(proposition)로 부르면 안 된다.

\subsection{대수(algebra)·수치(numerics)·확률(probability) 층(layer)에 남은 의무}

frozen paper의 MINCORE 전선은 KeyGen \identifier{STOP-KG-NTT}, Sign
\identifier{STOP-SIGN-CHAL-MODE2}, Verify \FrozenPaperGate 로 정리되었다.
post-freeze에는 \identifier{rq_mul_coeff_foldr_to_bigi}, full-NTT spectral action,
native \(R_q\) 행곱과 HAETAE 표현 운송이 닫혔지만, actual ExpandVecA와 기존
security-side synthetic object의 의미론 불일치가 \CurrentGate 를 만든다.
이와 별개로 actual Verify의 2-by-4 NTT/CRT/freeze 절차 의미론도 frozen
의무로 남는다. 전체 구현 refinement에는 다음 독립적인 의무들이 남는다.

\begin{description}[leftmargin=39mm,style=nextline]
  \item[public-key algebra/NTT]
    \(R_q\)-모듈(module) 계산, NTT 표현, public-key 관계가 실제 구현과 일치함.
  \item[sampler distribution]
    uniform/small/hyperball sampler와 실제 random tape의 출력분포(output
    distribution)가 논문 분포(distribution)와 일치하거나 명시적 거리
    상계(distance bound)를 가짐.
  \item[rejection과 종료]
    KeyGen/Sign retry가 종료하고, 수락 분포(accepted distribution)와 retry loss가
    계산됨.
  \item[highbits/LSB]
    실제 coefficient decomposition, rounding, packing이 논문의
    정수 분해(integer decomposition)와
    일치하며 challenge 양쪽 입력이 같음.
  \item[FIPS202 transcript]
    모든 domain separator, 길이, absorb 순서, squeeze 길이가 고정 명세와
    byte-for-byte 일치함.
  \item[FFT tie policy]
    fixed-point FFT와 경계 동률 처리까지 analytic predicate와 연결됨.
  \item[two-transcript extraction]
    실제 byte-level challenge와 accepted signature instance를 paper/ROM reduction의
    두 transcript 추출에 연결함.
\end{description}

\subsection{Week 16 frozen KeyGen의 KG--NTT--MUL 종료 경계}

actual \identifier{_kp_m23_matrix}와 \identifier{_keypair_finalize_m23}의
투명한 two-call harness는 pre-finalization \(bp\) snapshot, finalizer semantics,
KG-2-like low/high 분해, adjusted \(s_2\), snapshot-only mod-\(2q\) 영 합동식을
같은 사후조건에 둔다
\cite{week16-mincore-plan,week16-kg-report,week16-kg-ntt-mul-report}. 이는
\claimid{OBL-MINCORE-KEYGEN}의 실제 coefficient 경계에서 의미 있는
\statuspartial 진전이다.

matrix helper의 결과는 \identifier{KeygenM23ArithmeticSpec.output_row}로
표현된다. 새 \theoremname{output_row_from_mode2_ntt_words}는 이것을 actual full
NTT, pointwise row product, inverse NTT와 Montgomery 표현으로 정확히 rewrite하며,
\theoremname{actual_m23_matrix_snapshot_rows_explicit}는 투명한 harness의 두 active
row에 그 표현을 적용한다.
frozen paper 시점에는 그 spectral 표현이 native \(R_q\) negacyclic 곱과 같다는
compiled 정리와, \identifier{Rq.poly}를 보안 모형의
\(A_{\mathrm{gen}}s_{\mathrm{gen}}\) 표현으로 옮기는 어댑터가 없었다.
KG-NTT-MUL 감사는 누락 leaf를 odd-root orthogonality/full-NTT convolution과
이 어댑터로 특정하고 \identifier{STOP-KG-NTT}로 종료했다. frozen 판정에서
다음 승격은 금지된다.

\begin{itemize}
  \item snapshot의 \(\mathit{pre\_bp}\)를 근거 없이
        \(A_{\mathrm{gen}}s_{\mathrm{gen}}\)으로 이름 바꾸기;
  \item KG-2-like coefficient 분해에서 faithful KG-1 또는 KG-3를 결론내기;
  \item adjusted \(s_2\) 한 성분에서 complete KG-4 vector를 결론내기;
  \item snapshot 합동식을 논문식 \(A s=qj\pmod{2q}\)로 승격하기;
  \item 이 two-call harness에서 sampler, retry, packing, public API 또는 KeyGen
        종료성을 추론하기.
\end{itemize}

\subsection{post-freeze KG 진전과 새로운 의미론 경계}

frozen 뒤의 shared NTT theory는 pure coefficient normalization, full-NTT
spectral action과 임의 열 수 행곱을 증명했다. Mode~2에서는 KeyGen 3열 및 Verify
4열의 순수 specialization이 모두 컴파일된다. KeyGen 쪽은 native \(R_q\) 행곱에서
\identifier{HAETAE_Algebra.poly_dot}까지 운송되었다
\cite{post-freeze-kg-rq-haetae-report}. 단, Verify 4열 정리는 actual helper
loop를 열지 않으므로 \FrozenPaperGate 를 해제하지 않는다.

다음 감사는 actual \identifier{avec}의 SHAKE128 framing, nonce 515/516,
little-endian candidate와 \(<64513\) rejection을 paper \(a\)에 연결하고 paper
\(qj=q(1,0)^T\)를 snapshot zero equation에 독립적으로 더했다. 그러나 다음 세
대상은 서로 같지 않다.

\begin{center}
actual SHAKE128/rejection \(a\)
\(\neq\) synthetic security generator,
\qquad paper \(qj\)는 별도 correction.
\end{center}

compiled counterexample는 security-named row~0 coefficient~0이 401이고 paper
\(qj[0,0]=64513\)임을 보인다. all-zero seed trace의 actual 값도 44985다.
따라서 \CurrentGate 의 수리는 synthetic
\identifier{public_rounding_vector_seed}를 actual expansion 의미론으로 교체 또는
정제하고 coefficientwise actual-array representation을 증명하는 일이다
\cite{post-freeze-kg-actual-avec-qj-report}. 이 경계를 넘기 전에는 다음 승격도
금지된다.

\begin{itemize}
  \item synthetic object를 actual \identifier{avec} 또는 paper \(qj\)로 재명명;
  \item 별도 \(qj\) correction을 modulo-\(q\) HAETAE vector 안에서 보존된다고 가정;
  \item 새 actual-carrier 전용 coefficient-product를 기존 security
        matrix/secret만으로 정의한 일반 ring \identifier{matrix_vec_mul} 정리 없이
        security-level width-6 mod-\(2q\) 식으로 재명명;
  \item deterministic partial correctness에서 sampler distribution, uniformity,
        losslessness 또는 termination을 추론.
\end{itemize}

\sourcepath{KgFaithfulAugmentedModel.ec}는 이 포장을 위한 actual width-6 carrier와
generated-product transport를 정의하고 standalone fresh compile을 통과한다.
\theoremname{actual_faithful_augmented_productE}는 전용 structured
coefficient-product를 head, generated HAETAE \identifier{poly_dot}, adjusted
error의 합으로 전개한다.
\theoremname{checked_mode2_parent_m23_finalize_faithful_augmented_paper_as_qj}가
actual checked parent에 대해 이 product와 faithful \(qj\)의 합동식을
coefficientwise 증명한다. product는 actual matrix/secret만을 입력으로 받으며,
중간 2-by-3 block은 명시된 열에서 복원된 \identifier{poly_dot}이다.
따라서 이 작업의 판정은 \FaithfulAugmentedVerdict 다. 그러나 이 전용
mod-\(2q\) evaluator는 기존 mod-\(q\) security \identifier{matrix_vec_mul}이 아니고,
security-side ExpandVecA equivalence도 아니다. 따라서 \CurrentGate 가 사라지지는
않는다.

\sourcepath{ImplementationSecurityGoals.ec}는 이 항목들을
\identifier{obl_*} operation으로 열거한다. 이 파일이 컴파일된다는 것은
목표의 이름과 합성 인터페이스가 well-formed하다는 뜻이지 각 boolean이
참임을 증명했다는 뜻이 아니다.

\subsection{최상위 기능정확성은 아직 명세다}

\sourcepath{TopLevelGoals.ec}는
\begin{align*}
  \mathsf{FC}_{\mathrm{KG}}
    &:\quad \mathsf{JKeyGen}_2=\mathsf{PaperKeyGen}_2,\\
  \mathsf{FC}_{\mathrm{Sign}}
    &:\quad \forall sk,M,
      \mathsf{JSign}_2(sk,M)=\mathsf{PaperSign}_2(sk,M),\\
  \mathsf{FC}_{\mathrm{Verify}}
    &:\quad \forall vk,M,\sigma,
      \mathsf{JVerify}_2(vk,M,\sigma)
      =\mathsf{PaperVerify}_2(vk,M,\sigma)
\end{align*}
를 boolean specification으로 둔다. 여기서 Sign/KeyGen의 등식(equality)은
확률분포(probability distribution)의 등식(equality in distribution)을 포함한다.
현재의 prefix, memory, transcript, codec 정리(theorem)는 이 목표에
필요한 결정론적 하위 사각형이지만 전체 등식 자체는 아니다.

\subsection{보안 합성식은 아직 조건부 인터페이스다}

구현과 논문 스킴 사이의 목표는
\[
  \Adv^{\mathrm{euf\mbox{-}cma}}_{\mathrm{Jasmin}}
  \le
  \Adv^{\mathrm{euf\mbox{-}cma}}_{\mathrm{paper}}
  +\delta_{\mathrm{KeyGen}}
  +q_{\mathrm{Sign}}\delta_{\mathrm{Sign}}
  +\delta_{\mathrm{Verify}}
  +\delta_{\mathrm{Encoding}}
\]
의 모양이다. 현재 파일은 손실 변수(loss variable)를 선언하고
부등식(inequality)의 interface를 정의하지만,
각 \(\delta\)를 모두 닫아 최종 정리(theorem)로 합성하지 않았다.

논문 수준 reduction interface는 다시 MLWE, BST-MSIS, ROM 및 rejection/fork
오차로 내려간다. \sourcepath{SecurityAssumptions.ec}는 최종적으로 남길 수 있는
암호학적 인터페이스를 MLWE, BST-MSIS, ROM 세 종류로 제한한다. 하지만 구체적
게임 instance와 매개변수화된 reduction이 구현의 정확한 byte-level scheme에
연결되었다는 뜻은 아니다.

\begin{caution}[선언 종류의 구분]
연산(operation), 가정(assumption), 정리(theorem)는 서로 다르다.
\identifier{op obl_sampler_distribution : bool.}처럼 본문 없이 선언된 operation은
목표를 표현하는 추상 기호(abstract symbol)다. 그것을 참이라고 가정하는 axiom과 동일하지
않고, 그것을 참이라고 증명한 theorem도 아니다. 프로젝트의 검증 스크립트는
authored axiom과 proof hole을 별도로 금지·검색한다.
\end{caution}

\subsection{신뢰 경계와 범위 밖 항목}

현재 증거 사슬은 다음을 신뢰하거나 고정한다.

\begin{itemize}
  \item SHA-256가 기록된 46쪽 HAETAE 명세 PDF \cite{haetae-spec};
  \item pinned \sourcepath{haetae-ref-jasmin/} 구현과 focused extraction
        \cite{pinned-jasmin};
  \item EasyCrypt kernel, Jasmin model, 호출된 SMT prover와 build toolchain;
  \item source hash와 generated extraction/certificate manifest.
\end{itemize}

다음은 이 sprint의 결론에 포함되지 않는다.

\begin{itemize}
  \item mode~2 이외의 parameter set;
  \item 물리적 side channel, fault, RNG 품질;
  \item memory safety나 constant-time이 EUF-CMA에서 자동으로 따라온다는 주장;
  \item 실제 하드웨어/컴파일러 전체의 end-to-end correctness;
  \item full KeyGen, Sign, Verify 또는 구현 EUF-CMA가 완료되었다는 주장.
\end{itemize}

\subsection{현재 주장 행렬}

\begin{longtable}{@{}L{47mm} L{27mm} L{70mm}@{}}
\toprule
주장 & 상태 & 정확한 경계 \\
\midrule
\endfirsthead
\toprule
주장 & 상태 & 정확한 경계 \\
\midrule
\endhead
packed SK의 VK prefix & \statusproved &
실제 generated packer/Internal KeyGen의 반환 시 첫 992바이트 \\
raw \(\mu\) generated seam & \statusproved &
raw branch, 관련 key/pre/message region에서 64-to-32 prefix \\
position-64 challenge suffix & \statusproved &
동일 초기 state/position에서 32-byte absorb \\
actual API key copy/address/frame & \statusproved &
각 명명된 valid/canonical/disjoint 계약 아래 \\
accepted Verify tail/input binding & \statusproved &
accept \(\Rightarrow\) tail/hash inputs 방향 \\
stored observed \(\mu\) equality & \statuspartial &
product/replay transport 부재 \\
signature prefix inverse & \statusproved &
canonical challenge/low, 앞 1056바이트 \\
HBZ leaf inverse & \statusproved &
앞 1024 coefficient, rANS 제외 \\
actual rANS table와 pure step & \statusproved &
concrete 13-symbol table, 명시된 state range \\
pure normalization byte trace & \statusproved &
canonical symbol list, actual loop 아님 \\
actual array/list·word-step bridge & \statusproved &
1024-symbol suffix 점화식(recurrence), literal table word 등식(equality) \\
actual final serialization leaf & \statusproved &
네 번의 실제 store와 누적 trace의 little-endian 직렬화(serialization) 합성 \\
actual encoder success size & \statusproved &
반환한다면 실패이거나 \(0\le off\le1020\), \(4\le1024-off\le1024\) \\
actual inner normalization & \statusproved &
byte segment·frame·no-underflow와 정확한 순수 \(0..2\) exit \\
actual encoder trace refinement & \statusproved &
성공 조건부 부분정확성(success-conditioned partial correctness)의 전체 suffix 등식(equality) \\
actual decoder trace refinement & \statusproved &
exact-trace 부분정확성(partial correctness), 1024-symbol 등식(equality), 정확한 byte 소비와 tail frame \\
actual rANS harness control & \statusproved &
실제 호출 순서와 성공/실패 분기만 완료 \\
actual rANS core inverse & \statusproved &
actual encoder/copy/decoder의 성공 조건부 부분정확성(success-conditioned partial correctness) \\
actual encoder success witness & \statusproved &
fixed all-zero/all-6 반환 실행의 failure exclusion; termination/losslessness 없음 \\
production full-HBZ wrapper inverse & \statusproved &
prepare/apply, rANS core와 production lift의 성공 조건부 합성 완료 \\
actual KeyGen matrix/finalize snapshot & \statuspartial &
KG-2-like 분해와 snapshot mod-\(2q\) identity는 증명;
frozen paper에서는 마지막 NTT-word rewrite 뒤 \identifier{STOP-KG-NTT} \\
post-freeze NTT/Rq/HAETAE 행곱 & \statusproved &
순수 spectral action, native \(R_q\) 3열/4열 행곱과 KeyGen 3열 HAETAE 표현 운송;
actual Verify procedure는 제외 \\
post-freeze actual \identifier{avec}/paper \(a,qj\) & \statuspartial &
결정론적 expansion과 별도 \(qj\) snapshot 합성은 증명;
synthetic security 의미론 불일치에서 \CurrentGate \\
faithful augmented coefficient-product 합성 & \statusproved &
standalone actual-carrier structured product \(\equiv qj\); frozen 82-target 밖이며
generic security-model \identifier{matrix_vec_mul}/ExpandVecA equivalence는 제외 \\
actual Sign accepted core & \statuspartial &
실제 세 helper의 호출·accept branch control만;
\identifier{STOP-SIGN-CHAL-MODE2} \\
actual Verify core sequence & \statuspartial &
실제 다섯 helper 순서, V-1/V-2/V-5/V-6 word projection, norm/tail local 결과;
frozen \FrozenPaperGate \\
\(h\) suffix codec & \statusdeferred &
Week~16 MINCORE KeyGen 범위 전환; inverse 미시도 \\
full signature codec & \statusblocked &
HBZ wrapper 뒤에도 \(h\), metadata, padding과 residual frame 잔여 \\
Sign output tail reach & \statuspartial &
prefix와 success-conditioned HBZ wrapper만 닫힘; \(h\), metadata/padding,
norm/hint/arithmetic 잔여 \\
FC-KeyGen/Sign/Verify & \statusspecified / \statuspartial &
KeyGen snapshot 포함 하위 결정론적 seam만 일부 증명 \\
구현 EUF-CMA & \statusspecified &
loss와 paper reduction 연결 미완 \\
\bottomrule
\end{longtable}
